\clearpage\fancyhead[L]{\small COMET.}\hypertarget{native-343-lower}{}\label{native-343-lower}
\begin{center}\textbf{COMET.}\par\vspace{4pt}\hrule\vspace{4pt}\fontsize{10}{12}\selectfont Monatliche Correspondenz. Vol. XXIV, pp. 406\,..\,411, October 1811.\par\vspace{4pt}\hrule\vspace{4pt}\end{center}

\par\indent In the previous issue we expressed the hope that in this installment we would be able to tell our readers something definite about the comet’s orbital period. That hope has gone unfulfilled. At first it seemed likely that the large and steadily growing discrepancy between the observations and the parabolic elements arose from an actual departure of the orbit from a parabola. Later, however, it turned out that a small correction to the elements was quite enough to bring the two into satisfactory agreement. Since it is of some interest to see how successive changes bring the elements into ever closer agreement with the sky, we begin today by presenting a table comparing all observations made up to 11 October with \textsc{Gauss}’s first revised elements (Mon. Corresp., September issue, p. 305) [p. 335 of this volume.].

\Needspace{41\baselineskip}\SourceLeaf{344}
Comparison of all observations of the comet with \textsc{Gauss}’s parabolic elements.

\begin{center}

\fontsize{10}{12}\selectfont\setlength{\tabcolsep}{7pt}\renewcommand{\arraystretch}{1.05}\begin{tabular}{ll|c|c|l}
\multicolumn{2}{c|}{1811} & \multicolumn{2}{c|}{Deviation} & Observer’s \\
\multicolumn{2}{c|}{Date of observation} & in R.A. & in decl. & name \\\hline
August & 22 & $+61''$ & $-57''$ & Bessel \\
 & 23 & $+70$ & $-79$ & Bessel \\
 & 23 & $+119$ & $-21$ & Olbers \\
 & 25 & $+116$ & $-101$ & Olbers \\
 & 26 & $+98$ & $-84$ & Olbers \\
 & 27 & $+99$ & $-99$ & Olbers \\
 & 27 & $+63$ & $-50$ & Bessel \\
 & 28 & $+34$ & $-68$ & Bessel \\
 & 28 & $+46$ & $-124$ & Olbers \\
 & 29 & $+128$ & $-26$ & Oriani \\
 & 30 & $+9$ & $-41$ & Olbers \\
 & 30 & $+43$ & $-55$ & Bessel \\
 & 31 & $+37$ & $-49$ & Oriani \\
 & 31 & $+47$ & $-120$ & Olbers \\
Sept. & 1 & $+25$ & $-43$ & Oriani \\
 & 2 & $+19$ & $-74$ & Oriani \\
 & 3 & $+10$ & $-92$ & Oriani \\
 & 3 & $+23$ & $-67$ & Olbers \\
 & 4 & $+25$ & $-73$ & Oriani \\
 & 4 & $+51$ & $-129$ & Olbers \\
 & 4 & $+30$ & $-9$ & Gauss \\
 & 6 & $-43$ & $-89$ & Oriani \\
 & 6 & $-7$ & $-50$ & Lindenau \\
 & 6 & $+9$ & $-11$ & Gauss \\
 & 7 & $-26$ & $-74$ & Lindenau \\
 & 7 & $-17$ & $-80$ & Bessel \\
 & 7 & $-54$ & $+20$ & Gauss \\
 & 8 & $-54$ & $-45$ & Oriani \\
 & 8 & $-41$ & $-44$ & Lindenau \\
 & 9 & $-61$ & $-38$ & Oriani \\
 & 9 & $-62$ & $-7$ & Lindenau \\
 & 9 & $-65$ & $-64$ & Gauss \\
\end{tabular}
\end{center}
\Needspace{38\baselineskip}
\begin{center}

\fontsize{10}{12}\selectfont\setlength{\tabcolsep}{7pt}\renewcommand{\arraystretch}{1.05}\begin{tabular}{ll|c|c|l}
\multicolumn{2}{c|}{1811} & \multicolumn{2}{c|}{Deviation} & Observer’s \\
\multicolumn{2}{c|}{Date of observation} & in R.A. & in decl. & name \\\hline
Sept. & 10 & $-116''$ & $-27$ & Lindenau \\
 & 10 & $-54$ & $-71$ & Schubert \\
 & 11 & $-70$ & $-73$ & Oriani \\
 & 11 & $-135$ & $-28$ & Lindenau \\
 & 11 & $-82$ & $-106$ & Bessel. \\
 & 12 & $-120$ & $-43$ & Oriani \\
 & 12 & $-37$ & $-78$ & Bessel \\
 & 12 & $-86$ & $-53$ & Schubert \\
 & 13 & $-111$ & $-118$ & Oriani \\
 & 13 & $-130$ & $-34$ & Lindenau \\
 & 13 & $-119$ & $-74$ & Schubert \\
 & 14 & $-135$ & $-94$ & Oriani \\
 & 14 & $-146$ & $-34$ & Lindenau \\
 & 14 & $-140$ & .\;.\;. & Gauss \\
 & 14 & $-139$ & $-30$ & Schubert \\
 & 15 & $-154$ & $-51$ & Lindenau \\
 & 15 & $-195$ & $-76$ & Gauss \\
 & 15 & $-158$ & $-65$ & Schubert \\
 & 16 & $-225$ & $-54$ & Lindenau \\
 & 16 & $-219$ & $-30$ & Gauss \\
 & 16 & $-184$ & $-52$ & Schubert \\
 & 17 & $-172$ & $-43$ & Oriani \\
 & 17 & $-261$ & $-29$ & Lindenau \\
 & 18 & $-247$ & $-34$ & Lindenau \\
 & 19 & $-287$ & $-66$ & Lindenau \\
 & 22 & $-443$ & .\;.\;. & Bessel \\
 & 23 & $-420$ & .\;.\;. & Bessel \\
 & 26 & $-626$ & $+53$ & Lindenau \\
 & 30 & $-839$ & $+130$ & Lindenau \\
Oct. & 2 & $-879$ & $+118$ & Gauss \\
 & 5 & $-1123$ & $+336$ & Lindenau \\
 & 11 & $-1271$ & $+626$ & Lindenau \\
\end{tabular}

\end{center}

\Needspace{17\baselineskip}\par\indent Professor \textsc{Gauss}, who gave us the great pleasure of spending 14 days at our observatory in the last days of September and at the beginning of October, was much inclined to see a trace of an elliptical or hyperbolic orbit in these discrepancies. While still here he began some calculations towards determining the orbit, but reached no definite result on the matter. On the second day after his return to Göttingen, however, he already sent us the interesting results he had obtained. We give them here in his own words. “Dated Göttingen, 14 October. I checked my latest calculations on the comet at once and soon found, as I had suspected, a slip of 10 minutes in writing them down. Once I corrected it the result was quite different, namely,

\begin{center}
\emph{that there is still no certain trace of any ellipticity in the orbit.}
\end{center}

\noindent My revised parabolic elements are as follows.

\Needspace{12\baselineskip}
\begin{center}\fontsize{10}{12}\selectfont\setlength{\tabcolsep}{4pt}\renewcommand{\arraystretch}{1.15}\begin{tabular}{l}
Passage through perihelion, 1811 Sept. 12\\
\quad $5^{\mathrm u}\,32^{\mathrm m}\,46^{\mathrm s}$ mean time in Göttingen\\
\end{tabular}\end{center}
\begin{center}\fontsize{10}{12}\selectfont\setlength{\tabcolsep}{3pt}\renewcommand{\arraystretch}{1.15}\begin{tabular}{llr}
Longitude of perihelion & .\;.\;.\;.\;.\;.\;. & $75^\circ\,4'\,43''$\\
Log. of the distance at perihelion & .\;.\;. & $0.015530$\\
Longitude of the ascending node & . & $140^\circ\,21'\,40''$\\
Inclination of the orbit & .\;.\;.\;.\;.\;. & $73\;4\;18$\\
\end{tabular}\end{center}
\noindent Motion retrograde

\Needspace{17\baselineskip}\SourceLeaf{345}
\par\indent I find the following constants for the coordinates referred to the equator.

\[
\begin{array}{ll}
x = \dfrac{\alpha.\sin(v+348^\circ\,50'\,38'')}{\cos\frac{1}{2}v} & \log \alpha = 9.91435 \\[2ex]
y = \dfrac{\ObservedBeta.\sin(v+171^\circ\,56'\,51'')}{\cos\frac{1}{2}v} & \log \ObservedBeta = 9.80148 \\[2ex]
z = \dfrac{\gamma.\sin(v+80^\circ\;0'\;3'')}{\cos\frac{1}{2}v} & \log \gamma = 0.01538 \\
\end{array}
\]

\par\indent These elements, incidentally, were calculated only hastily, and they could easily be made to fit the observations even better. I did not think it worth the trouble to do so just yet, as they suffice for comparison with the observations for some time to come and show that it is still impossible to say anything certain about whether the orbit is elliptical. I do not believe the orbital period can be less than 1000 years if the orbit is elliptical. I have just compared these elements with your observations and mine, which confirms what I said a moment ago.

\Needspace{25\baselineskip}
\begin{center}\fontsize{10}{12}\selectfont\setlength{\tabcolsep}{6pt}\renewcommand{\arraystretch}{1.15}\begin{tabular}{l|c|c|c|c}
 & \multicolumn{4}{c}{Difference}\\
 & \multicolumn{2}{c|}{in right ascension} & \multicolumn{2}{c}{in declination} \\
1811 & v. L. & G. & v. L. & G. \\\hline
Sept. 4 & --- & $+5''$ & --- & $+23''$ \\
6 & --- & $-7$ & --- & $+28$ \\
7 & --- & $-58$ & --- & $+61$ \\
9 & --- & $-38$ & --- & $-18$ \\
10 & $-69''$ & & $+26''$ & \\
11 & $-70$ & & $+39$ & \\
14 & $-12$ & $-9$ & $+26$ & \\
15 & $+5$ & $-17$ & $+7$ & $-18$ \\
16 & $-37$ & $-36$ & $+10$ & $+35$ \\
17 & $-43$ & & $+33$ & \\
18 & $+4$ & & $+16$ & \\
19 & $-27$ & & $+1$ & \\
26 & $-47$ & & $+63$ & \\
30 & --- & $-50$ & $+76$ & \\
October 2 & --- & $-29$ & & $+19$ \\
\end{tabular}
\end{center}

\par\indent The elements thus give right ascensions about half a minute too small --- apparently throughout the entire period of these observations --- and declinations about as much too large. This discrepancy could easily be removed without departing from the parabola or lessening the agreement with von \textsc{Zach}’s observations. In a letter I found here, \textsc{Olbers} has given me a brief account of his ingenious hypothesis about the comet’s tail. He assumes a substance produced by the comet and repelled by both the comet and the Sun. It accumulates where the two repulsive forces maintain a kind of equilibrium and forms a sort of hollow parabolic conoid. \textsc{Olbers} will write to you about it in more detail.
